% ch2_radar_camera_calibration.tex — Radar-Camera Sensor Registration and Calibration Methods
\section{Radar-Camera Sensor Registration and Calibration Methods}\label{sec:calibration}

The preceding chapter established that maritime object detection datasets remain orders of magnitude smaller than their automotive counterparts, and that manual annotation constitutes the principal bottleneck to scaling these collections. A radar-guided labeling pipeline---one that projects radar detections into camera imagery to generate bounding box annotations automatically---offers a compelling path toward closing this gap. However, the fidelity of any such cross-modal labeling system is bounded above by the accuracy of the geometric registration between the radar and camera coordinate frames. An error of even a few pixels in the projected radar centroid propagates directly into label noise, degrading downstream detector training. Precise sensor registration is therefore not merely a desirable system property but a strict prerequisite for the entire radar-guided labeling paradigm.

This chapter surveys the methods by which radar and camera sensors are brought into geometric correspondence. The discussion begins with the mathematical framework that underlies all extrinsic calibration formulations, proceeds through the three dominant families of calibration techniques---target-based, targetless, and motion-based---and then addresses the complementary problems of intrinsic camera calibration and temporal synchronization. Throughout, attention is drawn to the gap between the well-studied 77\,GHz automotive radar domain and the largely unaddressed X-band marine radar setting that motivates the present work.

%% =========================================================================
\subsection{Mathematical Framework for Extrinsic Calibration}\label{sec:math_framework}
%% =========================================================================

The spatial relationship between any two rigidly mounted sensors is described by six degrees of freedom (6DOF): three translational components and three rotational components. In the context of radar-camera calibration, the objective is to recover the rigid-body transformation that maps points expressed in the radar coordinate frame into the camera coordinate frame \cite{persic2019extrinsic}. This transformation is parameterized by a $3\times3$ rotation matrix $\mathbf{R} \in SO(3)$ and a $3\times1$ translation vector $\mathbf{t} \in \mathbb{R}^{3}$.

\subsubsection{Rotation Representation}

The rotation matrix $\mathbf{R}$ is an orthonormal matrix satisfying $\mathbf{R}^{\top}\mathbf{R} = \mathbf{I}$ and $\det(\mathbf{R}) = +1$, encoding the orientation of the radar frame with respect to the camera frame. Explicitly, $\mathbf{R}$ can be decomposed into three successive rotations about the coordinate axes, parameterized by the Euler angles $(\phi, \theta, \psi)$ corresponding to roll, pitch, and yaw, respectively:
\begin{equation}\label{eq:rotation_matrix}
  \mathbf{R} = \mathbf{R}_{z}(\psi)\,\mathbf{R}_{y}(\theta)\,\mathbf{R}_{x}(\phi),
\end{equation}
where
\begin{equation}
\mathbf{R}_{x}(\phi)=\begin{bmatrix}1&0&0\\0&\cos\phi&-\sin\phi\\0&\sin\phi&\cos\phi\end{bmatrix},\quad
\mathbf{R}_{y}(\theta)=\begin{bmatrix}\cos\theta&0&\sin\theta\\0&1&0\\-\sin\theta&0&\cos\theta\end{bmatrix},\quad
\mathbf{R}_{z}(\psi)=\begin{bmatrix}\cos\psi&-\sin\psi&0\\\sin\psi&\cos\psi&0\\0&0&1\end{bmatrix}.
\end{equation}
Alternative rotation parameterizations---axis-angle, unit quaternions, and rotation vectors---are often preferred in optimization to avoid gimbal lock and to simplify constraint enforcement \cite{peretroukhin2020motion}. In particular, the unit quaternion representation $\mathbf{q} = (q_w, q_x, q_y, q_z)$ with $\|\mathbf{q}\| = 1$ is widely used in hand-eye calibration formulations \cite{daniilidis1999handeye}.

\subsubsection{Homogeneous Transformation}

The rotation and translation are combined into a single $4\times4$ homogeneous transformation matrix:
\begin{equation}\label{eq:homogeneous}
  \mathbf{T}_{c}^{r} =
  \begin{bmatrix}
    \mathbf{R} & \mathbf{t} \\
    \mathbf{0}^{\top} & 1
  \end{bmatrix}
  \in SE(3),
\end{equation}
where the superscript and subscript denote the source (radar) and target (camera) frames. A point $\mathbf{P}_{r} = [x_r,\, y_r,\, z_r,\, 1]^{\top}$ expressed in homogeneous radar coordinates is transformed into the camera frame as:
\begin{equation}\label{eq:point_transform}
  \mathbf{P}_{c} = \mathbf{T}_{c}^{r}\,\mathbf{P}_{r} = \mathbf{R}\,\mathbf{P}_{r}^{(3)} + \mathbf{t},
\end{equation}
where $\mathbf{P}_{r}^{(3)}$ denotes the inhomogeneous (three-element) coordinate vector. Once transformed into the camera frame, the point is projected onto the image plane via the camera intrinsic matrix $\mathbf{K}$:
\begin{equation}\label{eq:projection}
  s\begin{bmatrix}u\\v\\1\end{bmatrix}
  = \mathbf{K}\,[\mathbf{R} \mid \mathbf{t}]\,\mathbf{P}_{r},
  \qquad
  \mathbf{K} = \begin{bmatrix}f_x & 0 & c_x \\ 0 & f_y & c_y \\ 0 & 0 & 1\end{bmatrix},
\end{equation}
where $s$ is a scalar depth factor, $(f_x, f_y)$ are the focal lengths in pixel units, and $(c_x, c_y)$ is the principal point. This projection model forms the basis of reprojection error metrics used throughout calibration optimization \cite{zhang2000flexible}.

\subsubsection{Implications for Marine Radar}

For automotive 77\,GHz radar, each detection typically provides range, azimuth, and Doppler velocity, but elevation information may be absent or coarsely resolved \cite{patole20174d}. Marine X-band radar presents an even more constrained geometry: the mechanically rotating antenna produces a plan-position indicator (PPI) display that records range and azimuth only, with no elevation whatsoever \cite{Skolnik2008}. Consequently, the transformation in Equation~\eqref{eq:point_transform} is underdetermined unless an elevation assumption is imposed---typically that all targets lie at sea level---reducing the 3D problem to a 2D-to-2D homography or a constrained 3D projection with a fixed height prior.

%% =========================================================================
\subsection{Target-Based Calibration Methods}\label{sec:target_based}
%% =========================================================================

Target-based calibration methods rely on purpose-built calibration artifacts whose positions and reflective properties are known a priori. These approaches yield the highest absolute accuracy but require controlled experimental conditions and cooperative placement of calibration objects.

\subsubsection{Corner Reflectors and Radar-Visible Targets}

Trihedral corner reflectors are the standard calibration target for radar systems. Their large and predictable radar cross section (RCS) ensures reliable detection at known positions, providing precise range-azimuth anchor points in the radar frame \cite{Skolnik2008}. When co-located with a visually identifiable marker---a painted panel, an LED array, or a distinctive geometric pattern---the same physical location can be identified in both sensor modalities, yielding a set of correspondence pairs for extrinsic parameter estimation.

Per\v{s}i\'{c} et al.\ \cite{persic2019extrinsic} developed a multi-target calibration system for a 77\,GHz radar--LiDAR--camera suite that exploits the RCS characteristics of corner reflectors to resolve the data association problem. By deploying multiple reflectors at varying ranges and azimuths, the authors obtained sufficient correspondences to solve for the full 6DOF extrinsic transformation using nonlinear least-squares optimization. The reported reprojection accuracy was 2.2 pixels RMSE in the horizontal direction and 3.0 pixels RMSE in the vertical direction, demonstrating that sub-3-pixel accuracy is attainable with careful target placement and sufficient geometric diversity among the calibration points.

\subsubsection{Homography-Based Approaches}

Sugimoto et al.\ \cite{sugimoto2004obstacle} proposed an early radar-camera registration method based on planar homography estimation. Under the assumption that both radar detections and camera observations lie on a common ground plane, the relationship between radar coordinates $(r, \theta)$ and pixel coordinates $(u, v)$ reduces to a $3\times3$ homography matrix $\mathbf{H}$:
\begin{equation}\label{eq:homography}
  s\begin{bmatrix}u\\v\\1\end{bmatrix}
  = \mathbf{H}\begin{bmatrix}x_r\\y_r\\1\end{bmatrix},
\end{equation}
where $(x_r, y_r)$ are the Cartesian radar coordinates on the ground plane. This formulation requires a minimum of four non-collinear point correspondences and is estimated robustly using RANSAC \cite{fischler1981random}. While the planar assumption is well-suited to flat roadway environments, it also applies naturally to the maritime setting where targets occupy a nominally planar sea surface. The principal limitation is that elevated structures---ship superstructures, masts, bridge decks---violate the planarity constraint and introduce systematic projection errors at close range.

\subsubsection{Checkerboard Patterns for Camera Intrinsics}

Target-based calibration of the camera intrinsic parameters is a prerequisite for accurate extrinsic estimation. Zhang's flexible calibration technique \cite{zhang2000flexible} uses images of a planar checkerboard pattern captured at multiple orientations to solve for the focal lengths, principal point, and lens distortion coefficients. The method minimizes the reprojection error between detected corner locations and their model-predicted positions via Levenberg--Marquardt optimization. This technique has become the de facto standard in both academic research and industrial practice, and is implemented in widely used libraries such as OpenCV. In the context of radar-camera systems, camera intrinsic calibration is performed independently and its results are treated as fixed during the subsequent extrinsic calibration stage \cite{Caesar2020nuScenes, Geiger2012KITTI}.

%% =========================================================================
\subsection{Targetless Calibration Methods}\label{sec:targetless}
%% =========================================================================

Target-based methods, while accurate, require physical access to the sensing environment and cooperative placement of calibration artifacts. These requirements are impractical for deployed maritime systems where sensors are permanently mounted on vessels or coastal infrastructure. Targetless calibration methods instead exploit naturally occurring correspondences in the sensor data streams.

\subsubsection{Trajectory Alignment}

When both sensors independently track moving objects over time, the resulting trajectories can be aligned to recover the inter-sensor transformation. The Kabsch algorithm \cite{kabsch1976solution} solves for the optimal rotation between two sets of corresponding 3D points by minimizing the root-mean-square deviation. Its extension by Umeyama \cite{umeyama1991least} additionally recovers a uniform scale factor, yielding a similarity transformation. In the radar-camera context, radar tracks provide range-azimuth trajectories while camera tracks provide pixel-domain trajectories. Under a known camera model, the pixel trajectories can be back-projected into bearing-only rays, and the correspondence between radar and camera tracks is established through temporal co-occurrence and motion consistency.

Wang et al.\ \cite{wang2023online} formalized this approach as an online track-to-track association framework for automotive radar-camera calibration. By continuously matching radar detections to camera-tracked objects using a joint cost function over spatial proximity and velocity consistency, the authors accumulate correspondences over time and solve for the extrinsic parameters using a Perspective-n-Point (PnP) formulation \cite{lepetit2009epnp}. The online formulation permits continuous recalibration, compensating for mechanical drift and thermal deformation---properties that are particularly relevant for shipboard installations subject to vibration and thermal cycling.

\subsubsection{Mutual Information and Feature-Based Methods}

Lin et al.\ \cite{lin2024targetless} proposed a targetless radar-camera calibration method that maximizes the mutual information between radar measurements and camera features. By treating the calibration problem as an information-theoretic alignment task, the method avoids explicit feature detection or correspondence establishment. Instead, it searches over the space of extrinsic parameters to find the transformation that maximizes the statistical dependence between the projected radar data and the image content. This approach is particularly robust to the sparse and noisy nature of radar point clouds, which often frustrate traditional feature matching techniques.

\subsubsection{Deep Learning Approaches}

The application of deep neural networks to sensor calibration has grown rapidly in recent years. Schneider et al.\ \cite{schneider2017regnet} introduced RegNet, a convolutional neural network that predicts the extrinsic calibration between a LiDAR and a camera from a single pair of depth and color images. The network is trained on synthetically perturbed calibrations and learns to regress the six extrinsic parameters directly. While originally developed for LiDAR-camera pairs, the RegNet architecture has been adapted to radar-camera calibration by substituting radar-derived depth or occupancy representations for the LiDAR depth image \cite{peng2024camera}.

More recently, end-to-end calibration networks such as RC-AutoCalib \cite{arXiv2025rcautocalib} have been proposed specifically for the radar-camera modality pair. These networks ingest raw radar tensors and camera images and output the extrinsic transformation without requiring intermediate detection or tracking stages. A comprehensive review by Peng et al.\ \cite{peng2024camera} categorizes deep learning calibration methods along axes of supervision level (fully supervised, self-supervised, unsupervised), input representation (point cloud, bird's-eye view, range-Doppler), and output parameterization (Euler angles, quaternions, direct matrix). The review observes that deep learning calibration methods typically achieve accuracy within 1--5 pixels of target-based methods on automotive benchmarks, but require large volumes of diverse training data that are not yet available for the maritime domain.

%% =========================================================================
\subsection{Motion-Based (Ego-Motion) Calibration}\label{sec:motion_based}
%% =========================================================================

Motion-based calibration exploits the ego-motion of the sensor platform itself rather than the motion of external targets. By observing the same platform motion from two different sensor viewpoints, the relative transformation between the sensors can be inferred.

\subsubsection{The Hand-Eye Calibration Formulation}

The classical formulation of motion-based calibration originates in robotics, where the problem of determining the transformation between a camera mounted on a robot end-effector (the ``eye'') and the end-effector itself (the ``hand'') was first posed by Shiu and Ahmad and subsequently solved in closed form by Tsai and Lenz \cite{tsai1989technique}. The hand-eye equation takes the form:
\begin{equation}\label{eq:handeye}
  \mathbf{A}\mathbf{X} = \mathbf{X}\mathbf{B},
\end{equation}
where $\mathbf{A}$ and $\mathbf{B}$ are the motions observed by the two sensors between successive time steps, and $\mathbf{X}$ is the unknown inter-sensor transformation. Given $n$ motion pairs $\{(\mathbf{A}_i, \mathbf{B}_i)\}_{i=1}^{n}$, the system is overdetermined for $n \geq 2$ (with non-parallel rotation axes) and can be solved via least-squares methods.

\subsubsection{Dual Quaternion Solution}

Daniilidis \cite{daniilidis1999handeye} reformulated the hand-eye problem using dual quaternions, which encode rotation and translation in a unified algebraic structure. The dual quaternion representation enables a simultaneous, closed-form solution for both $\mathbf{R}$ and $\mathbf{t}$, avoiding the sequential rotation-then-translation estimation of earlier methods. The resulting algorithm is computationally efficient, requiring only a singular value decomposition (SVD) of a $6n \times 8$ matrix, and is provably optimal under isotropic Gaussian noise. This formulation has been adopted in several subsequent radar-camera calibration systems \cite{persic2019extrinsic} and remains the preferred closed-form solution when sufficient motion diversity is available.

\subsubsection{Continuous-Time Batch Estimation}

A limitation of the classical hand-eye formulation is that it assumes discrete, synchronized motion observations from both sensors. In practice, radar and camera operate at different frame rates and are rarely hardware-synchronized. Furgale et al.\ \cite{furgale2013unified} addressed this challenge by proposing a continuous-time batch estimation framework in which sensor trajectories are represented as temporal basis spline (B-spline) curves. The continuous-time representation decouples the calibration problem from sensor frame rates: observations from asynchronous sensors are related through the common spline representation, and both temporal offsets and spatial extrinsics are estimated jointly. This approach is particularly relevant for maritime systems where the radar antenna rotation period (typically 1--3 seconds for X-band marine radars) is substantially longer than the camera frame interval, creating large temporal gaps between radar and camera observations of the same scene.

%% =========================================================================
\subsection{Intrinsic Camera Calibration}\label{sec:intrinsic_calib}
%% =========================================================================

Accurate intrinsic calibration is a prerequisite for meaningful extrinsic estimation, as errors in the camera model propagate directly into the projection of radar coordinates onto the image plane. Zhang's method \cite{zhang2000flexible} recovers the intrinsic parameter vector $\boldsymbol{\xi} = (f_x, f_y, c_x, c_y, k_1, k_2, p_1, p_2, \ldots)$, where $k_i$ are radial distortion coefficients and $p_i$ are tangential distortion coefficients. The calibration procedure captures $m$ images of a planar checkerboard at varied orientations and minimizes the total reprojection error:
\begin{equation}\label{eq:reprojection}
  \min_{\boldsymbol{\xi},\,\{\mathbf{R}_j,\,\mathbf{t}_j\}_{j=1}^{m}}
  \sum_{j=1}^{m}\sum_{i=1}^{n}
  \left\| \mathbf{u}_{ij} - \hat{\mathbf{u}}(\boldsymbol{\xi},\,\mathbf{R}_j,\,\mathbf{t}_j,\,\mathbf{X}_i) \right\|^{2},
\end{equation}
where $\mathbf{u}_{ij}$ is the observed pixel location of the $i$-th checkerboard corner in the $j$-th image and $\hat{\mathbf{u}}(\cdot)$ is its predicted location under the camera model. The optimization is initialized via a closed-form homography decomposition and refined iteratively.

For wide field-of-view cameras commonly used in maritime surveillance, lens distortion can exceed 20 pixels at the image periphery. Higher-order distortion models (e.g., the rational model with six coefficients or fisheye models) may be necessary to achieve sub-pixel reprojection accuracy across the full field of view \cite{zhang2000flexible}. Failure to adequately model distortion introduces spatially varying bias in the projected radar coordinates, which is particularly problematic for targets at extreme bearing angles near the edges of the camera's field of view.

%% =========================================================================
\subsection{LiDAR-Camera Calibration: Transferable Concepts}\label{sec:lidar_camera}
%% =========================================================================

The LiDAR-camera calibration literature is substantially more mature than its radar-camera counterpart, and many techniques developed for LiDAR-camera registration are directly transferable or adaptable to the radar-camera problem. Li et al.\ \cite{li2023targetless} provide a comprehensive survey that organizes LiDAR-camera calibration methods into four categories: target-based, feature-based, motion-based, and learning-based. This taxonomy maps cleanly onto the radar-camera methods discussed above.

Target-based LiDAR-camera methods typically use planar boards or polygonal cutouts that produce strong edges in both the point cloud and the image \cite{Geiger2012KITTI, zheng2021novel}. The analogous radar target is the corner reflector, though its point-like radar signature provides less geometric structure than a planar board in a LiDAR scan. Feature-based LiDAR-camera methods exploit edge correspondences, line features, or statistical correlations between depth discontinuities and image gradients; these techniques assume dense point clouds that are not available from radar. Motion-based and learning-based methods are more readily transferable because they operate on trajectory-level or representation-level abstractions that are less sensitive to point cloud density.

The KITTI benchmark \cite{Geiger2012KITTI} and nuScenes dataset \cite{Caesar2020nuScenes} have served as standard evaluation platforms for LiDAR-camera calibration. Both provide factory-calibrated extrinsic parameters as ground truth, enabling quantitative comparison of calibration algorithms. No analogous benchmark exists for radar-camera calibration in either the automotive or maritime domain, complicating the evaluation of proposed methods.

%% =========================================================================
\subsection{Temporal Synchronization Challenges}\label{sec:time_sync}
%% =========================================================================

Even a geometrically perfect extrinsic calibration degrades rapidly if the two sensor data streams are not temporally aligned. At maritime vessel speeds of 5--15 knots and typical ranges of 0.5--5\,km, a timing offset of 100\,ms produces a target displacement of 0.26--0.77\,m, which may translate into several pixels of projection error depending on range and camera resolution.

\subsubsection{Asynchronous Sensor Architectures}

Radar and camera systems rarely share a common clock. Automotive 77\,GHz radars typically output detections at 10--20\,Hz, while cameras operate at 15--30\,fps. In the nuScenes system \cite{Caesar2020nuScenes}, hardware triggering and a central timing bus achieve sub-millisecond synchronization across all sensors. Such infrastructure is unavailable in typical maritime installations, where the radar interface provides PPI sweeps at the antenna rotation rate (20--48\,rpm for X-band marine radars) and the camera free-runs at its own frame rate \cite{Skolnik2008}.

\subsubsection{The Rotating Antenna Problem}

A challenge unique to mechanically scanned marine radar is that the antenna sweeps through $360^{\circ}$ over a rotation period of 1.25--3\,s. Unlike a phased-array automotive radar that illuminates its entire field of view quasi-simultaneously, a rotating marine radar observes different azimuth sectors at different instants within each sweep. A camera frame captured at time $t$ corresponds to a single azimuth sector of the radar sweep, not to the full PPI image. Naive association of a complete PPI scan with a single camera frame introduces bearing-dependent timing errors that worsen for targets at azimuths far from the instantaneous antenna pointing direction at time $t$.

Li and Pollefeys \cite{li2011temporal} studied temporal calibration for asynchronous sensor pairs and demonstrated that temporal offsets can be estimated jointly with spatial extrinsics by exploiting motion consistency constraints. The continuous-time framework of Furgale et al.\ \cite{furgale2013unified} extends this idea by representing sensor trajectories as continuous functions of time, enabling interpolation between discrete observations. For maritime radar-camera systems, such continuous-time methods offer a principled solution to the rotating antenna synchronization problem, provided that the antenna angle encoder signal is available to timestamp individual radar returns within each sweep.

%% =========================================================================
\subsection{Maritime-Specific Challenges}\label{sec:maritime_challenges}
%% =========================================================================

The vast majority of radar-camera calibration research targets the automotive domain, where 77\,GHz millimeter-wave radar is the standard sensor. Transferring these methods to the maritime setting introduces several challenges arising from fundamental differences in sensor characteristics, platform dynamics, and operational environment.

\subsubsection{X-Band Marine Radar Versus 77\,GHz Automotive Radar}

Table~\ref{tab:radar_camera_comparison} summarizes the key differences between X-band marine radar and 77\,GHz automotive radar that affect calibration methodology and achievable accuracy.

\begin{table}[htbp]
  \centering
  \caption{Comparison of X-band marine radar and 77\,GHz automotive radar characteristics relevant to calibration.}
  \label{tab:radar_camera_comparison}
  \begin{tabular}{@{}lll@{}}
    \toprule
    \textbf{Parameter} & \textbf{X-Band Marine Radar} & \textbf{77\,GHz Automotive Radar} \\
    \midrule
    Operating frequency      & 9.3--9.5\,GHz                 & 76--81\,GHz \\
    Wavelength               & $\sim$32\,mm                   & $\sim$3.9\,mm \\
    Maximum range            & 50--130\,km                    & 0.1--0.3\,km \\
    Range resolution         & 5--50\,m                       & 0.15--0.5\,m \\
    Azimuth beamwidth        & $1.2^{\circ}$--$5.2^{\circ}$  & $1^{\circ}$--$15^{\circ}$ \\
    Elevation information    & None (2D PPI)                  & Limited or none \\
    Scan mechanism           & Mechanical rotation            & Electronic beam steering \\
    Scan rate                & 20--48\,rpm                    & 10--20\,Hz (full FOV) \\
    Output format            & PPI image (polar)              & Detection list (range, azimuth, Doppler) \\
    Doppler capability       & Typically none                 & Yes (FMCW) \\
    Typical mounting height  & 15--40\,m (ship mast)          & 0.3--0.8\,m (vehicle bumper) \\
    \bottomrule
  \end{tabular}
\end{table}

Several of these differences have direct consequences for calibration. The absence of elevation data in X-band PPI radar means that the full 3D-to-2D projection model (Equation~\eqref{eq:projection}) is underdetermined without an elevation prior. The coarse range resolution (meters to tens of meters) and azimuth resolution ($>1^{\circ}$) limit the precision of radar-side point localization, placing an intrinsic floor on achievable reprojection accuracy that is typically far above the sub-pixel precision attainable with camera corner detection \cite{Skolnik2008, Richards2022}. The PPI output format requires centroid extraction from extended radar returns, introducing additional localization uncertainty compared to the point-detection output of automotive FMCW radar \cite{patole20174d}.

\subsubsection{Platform Dynamics and Mounting Geometry}

Maritime platforms are subject to continuous six-degree-of-freedom motion---surge, sway, heave, roll, pitch, and yaw---driven by sea state. On vessels, this motion can exceed $\pm15^{\circ}$ in roll and $\pm5^{\circ}$ in pitch under moderate sea conditions \cite{Ward2006}. Such motion invalidates any static calibration unless the calibration is referenced to a stabilized frame or compensated in real time using inertial measurements.

Furthermore, the large physical separation between maritime radar and camera installations introduces substantial lever-arm effects. A marine radar antenna mounted on a mast 30\,m above the waterline and a camera mounted at bridge level 20\,m above the waterline are separated by a vertical baseline of 10\,m and potentially a horizontal offset of several meters. Platform rotation amplifies these offsets: a $5^{\circ}$ roll about the ship's longitudinal axis displaces the radar antenna tip by $30 \times \sin(5^{\circ}) \approx 2.6$\,m laterally, while the camera moves by $20 \times \sin(5^{\circ}) \approx 1.7$\,m. The differential displacement of approximately 0.9\,m introduces time-varying extrinsic errors that must be compensated.

Coastal surveillance installations on fixed infrastructure (piers, towers, buildings) avoid vessel motion but introduce their own challenges. The radar antenna is often mounted on a dedicated mast structure separated from the camera by tens of meters, and the installation geometry is dictated by structural constraints rather than calibration convenience. Han et al.\ \cite{han2020autonomous} describe a coastal radar-camera system for autonomous surface vehicle navigation in which the radar and camera are co-located on a fixed tower, but note that thermal expansion and wind-induced vibration can shift the extrinsic calibration by several pixels over the course of a day.

\subsubsection{Environmental Factors}

The maritime environment presents clutter and propagation conditions that have no automotive analogue. Sea clutter from wave returns can produce false radar detections that, if mistakenly associated with camera features, corrupt calibration correspondences \cite{Ward2006}. Multipath propagation over the sea surface creates ghost targets at incorrect ranges. Atmospheric ducting in the marine boundary layer bends radar beams, causing anomalous propagation that shifts the apparent azimuth and range of targets \cite{Skolnik2008}. On the camera side, sun glare reflected off the water surface can saturate image regions and prevent feature detection. These environmental factors collectively reduce the number of reliable sensor correspondences available for calibration and increase the importance of robust estimation techniques such as RANSAC \cite{fischler1981random}.

%% =========================================================================
\subsection{Summary and Research Gaps}\label{sec:calib_gaps}
%% =========================================================================

This chapter has surveyed the three principal families of radar-camera calibration methods---target-based, targetless, and motion-based---along with the supporting infrastructure of intrinsic camera calibration, temporal synchronization, and the transferable insights from LiDAR-camera calibration research. The mathematical framework for extrinsic calibration is well-established: the 6DOF rigid-body transformation in $SE(3)$ and its projection through the camera intrinsic model provide a sound theoretical basis for all methods reviewed.

Within the automotive domain, each calibration family has demonstrated practical viability. Target-based methods achieve sub-3-pixel reprojection accuracy \cite{persic2019extrinsic}, online targetless methods enable continuous recalibration \cite{wang2023online}, and continuous-time batch estimation methods handle asynchronous sensor architectures \cite{furgale2013unified}. Deep learning approaches offer the prospect of fully automatic calibration without specialized targets or motion requirements \cite{schneider2017regnet, arXiv2025rcautocalib}.

However, the literature reveals a pronounced gap at the intersection of these methods and the X-band marine radar domain. The key observations are as follows:

\begin{enumerate}[leftmargin=*]
  \item \textbf{No established X-band marine radar to camera calibration method exists.} The overwhelming majority of published radar-camera calibration work uses 77\,GHz automotive radar. The handful of maritime radar-camera studies \cite{han2020autonomous, xu2025multistage} address narrow operational scenarios and do not provide general-purpose calibration frameworks.

  \item \textbf{The 2D PPI output of marine radar lacks elevation information}, rendering standard 3D-to-2D projection models underdetermined. A principled treatment of the elevation ambiguity---whether through sea-surface planarity assumptions, homographic models, or probabilistic elevation priors---has not been established for X-band radar-camera systems.

  \item \textbf{Temporal synchronization between mechanically rotating radar and free-running cameras} presents a unique challenge that has no direct automotive analogue. The continuous-time estimation frameworks of Furgale et al.\ \cite{furgale2013unified} provide a theoretical solution, but their application to marine radar with per-spoke timestamping has not been demonstrated.

  \item \textbf{Platform dynamics on vessels} introduce time-varying extrinsic parameters that violate the static calibration assumption underlying most existing methods. While inertial compensation is conceptually straightforward, its integration with radar-camera calibration for maritime platforms has received limited attention.

  \item \textbf{No benchmark dataset with ground-truth calibration exists} for X-band marine radar-camera systems, precluding quantitative comparison of competing approaches. The KITTI \cite{Geiger2012KITTI} and nuScenes \cite{Caesar2020nuScenes} benchmarks serve this role for automotive LiDAR and radar, but no maritime equivalent is available.

  \item \textbf{Coarse spatial resolution of X-band radar} places an intrinsic lower bound on calibration accuracy that is substantially higher than the sub-pixel precision achievable with automotive radar. Characterizing this bound and understanding its implications for downstream labeling quality is an open problem.
\end{enumerate}

These gaps collectively indicate that developing a calibration methodology specifically tailored to X-band marine radar-camera systems is a necessary precondition for the radar-guided camera labeling pipeline that motivates this review. The following chapter addresses the next link in that pipeline: the projection of calibrated radar detections into camera pixel coordinates and the generation of bounding box annotations from those projections.
